\section{Dabbling in the Quantum World}

\begin{quotex}
Nothing appears to be more certain than our scientific knowledge of the physical universe. But then, what is the physical universe? We are told that it consists of space, time, and matter, or of space-time and energy, or perhaps of something else still more abstruse and even less imaginable; but in any case we are told in unequivocal terms what it excludes: as all of us have learnt, the physical universe is said to exclude just about everything which from the ordinary human point of view makes up the world. Thus it excludes the blueness of the sky and the roar of breaking waves, the fragrance of flowers and all the innumerable qualities---half-perceived and half-intuited---that lend color, charm, and meaning to our terrestrial and cosmic environment. In fact, it excludes everything that can be imagined or conceived, except in abstract mathematical terms.
\flright{\textsc{Wolfgang Smith}, \emph{Cosmos and Transcendence}}

Thus nature gets credit which should in truth be reserved for ourselves: the rose for its scent, the nightingale for his song, and the sun for its radiance. The poets are entirely mistaken. They should address their lyrics to themselves, and should turn them into odes of self-congratulation and the excellency of the human mind. Nature is a dull affair, soulless, scentless, colorless; merely the hurrying of material, endlessly, meaninglessly.
\flright{\textsc{Alfred North Whitehead}, \emph{Science and the Modern World}}

We are generally in agreement in thinking that humanity has arrived at an important turning point in its history. The Cartesian spirit which destroyed scholastic philosophy is now in turn being left behind. The logic of history demands a new spirit. The divorce between traditional knowledge, of which religion is a trustee, and acquired knowledge, the fruit of science, threatens to make sterile the Christian civilization which in origin is so rich with promises. Yet it is an aberration to believe that Science by its very nature is opposed to Tradition, and it must also be firmly stressed that Tradition does not include any tendency opposed to Science.
\flright{\textsc{Boris Mouravieff}, \emph{Gnosis}}
\end{quotex}


\paragraph{Quantum Physics}


Despite the incredible accuracy of its predictions and the elegance of its mathematical models, quantum mechanics still remains unintelligible. In other words, there is no answer to the ``what'' question (which is qualitative), but rather there are several wildly different interpretations -- all mathematically equivalent---of the equations. Wolfgang Smith offers his own interpretation in terms of Traditional metaphysics.

Metaphysics envisages different strata or degrees of being from the most dense to the most sublime. The naïve mind takes our everyday, corporeal, plane as the only reality. The \emph{materia prima} is the lowest level. According to \textbf{Rene Guenon} (see Chapter II of \emph{The Reign of Quantity}), this is the passive principle of universal manifestation or pure potentiality. As such, it is the ``subject of all the substantial transformations of the corporeal universe''. Hence, materia prima is unintelligible since things are known through their form. In the Vedanta, this level is called Prakriti or \emph{avyakta} in the Sankhya.

The next level is materia secunda, which is material, extended, quantitative, not yet a substance. That is, it has quantity but not quality and Guenon claims that this is what physicists actually mean by matter.

Unfortunately, modern science knows nothing about such levels and tries to understand things horizontally rather than vertically. Specifically, it view the world on a quantitative scale from subatomic particles, to atoms, to man, all the way to the largest structures of solar systems, galaxies, etc. The larger structures are not grasped as wholes, but are assumed to be built up from lower structures.

Now when this is applied to quantum events, we start seeing anomalies. For example, we consider the sun to be a ``ball'' and we know about balls in our corporeal world. Hence, we envision an electron, say, as just a very small ball. Yet, it doesn't always act as ball (i.e., a particle), but sometimes like a wave. Smith interprets that to mean that what we ``see'' are not electrons, but rather probabilities; i.e., electrons as such are quantitative, not qualitative.

Smith takes this to mean that the quantum world is actually a stratum of matter between the corporeal plane and the materia secunda. This, then, is the physical world, the world studied by physicists. In fact, we never actually experience the physical world directly, but we know it indirectly through its interactions with the corporeal world (i.e., instruments, cloud chambers, and other measuring devices).

Now these two strata are intimately related, since behind every object X in the corporeal world, there is a physical object SX in the subcorporeal world that corresponds to it. The physicist investigates SX, not X. This is Smith's unique contribution to the proper understanding of quantum mechanics.

The alternative horizontal view is that the object X is ``composed'' of various subatomic particles. Hence, we hear that the table we see is ``really'' mostly empty space. However, there is obviously no physical explanation for the table, since quantity cannot explain quality. That is hardly an increase in understanding, but rather a reversion to the first philosophers who believed that corporeal substances were just aggregations of atoms. It ignores the deepening of understanding from that view to Plato, Aristotle, and Thomas Aquinas.

\paragraph{Bell's Theorem}


In 1964, the physicist John Bell discovered a flaw in the equations of quantum mechanics. Up to that point, it was simply assumed that physical objects acted locally, hence were constrained to the known physical forces. In particular, signals between object could not exceed the speed of light. Bell's Theorem asserts, on the contrary, that reality is nonlocal. Not just in theory, this has been verified experimentally.

Smith sees this as a fundamental limitation of physics: it cannot in itself tell us about the universe, but it can only describe phenomena. It takes a metaphysical explanation to deal with the strata of cosmic reality. Smith relates this to the teaching on Heaven, Earth, and the metaxy as Plato called it, or the spiritual, corporeal, and the intermediate. Corporeal things are intermediary things subjected to spatio-temporal constraints. Smith claims that the alchemical processes of \emph{solve} and \emph{coagula} represent the transitions between those two strata.

Nonlocality can be understood by the traditional distinction between the corporeal and the intermediary levels; in the Vedanta, these are called the \emph{sthula} or gross level and the \emph{sukshma} or subtle level. Thus physics is consistent with this teaching:

\begin{itemize}


\item Quantum mechanics affirms the ontological discontinuity between the physical, subcorporeal, world and the corporeal by state vector collapse

\item Nonlocality illustrates the discontinuity between the corporeal and the intermediary
\end{itemize}


Physics qua physics cannot explain this. The physicist himself is at the corporeal level. On the one hand, he is looking downward toward the materia secunda, away from essences. Hence, he simply cannot see the corporeal and perforce the intermediary levels.

\paragraph{Summing It Up}


This has ramifications beyond just physics. Corporeal beings have a ``double birth'':

\begin{itemize}


\item A pre-temporal birth in the divine creative Act (vertical)

\item A temporal birth bringing it into the corporeal domain (horizontal)
\end{itemize}


Debates about science and religion are pointless without this understanding. Being blind to the vertical strata above the corporeal world, science will misinterpret metaphysical explanations as pertaining just to the corporeal world.

A couple of decades ago, quantum mechanics ``light'' was made popular in New Age circles in a few books. We don't need it to make us aware of higher dimensional reality, just the opposite, actually. I recall one woman, a school guidance counselor, who assured me that she was ``dabbling'' in quantum mechanics. I was the wrong person to hear that, since I could actually do the math involved, while she just wanted the false authority of science to back up her weird worldview. We hope that readers aren't just dabbling, but are reading with the aim of achieving an intellectual conversion.

Forthcoming posts will deal with Cosmology, Biology, and Psychology.


\flrightit{Posted on 2016-05-19 by Cologero}

\begin{center}* * *\end{center}

\begin{footnotesize}
\begin{sffamily}
\texttt{Sigurd on 2017-02-03 at 18:02 said:}

Hello, I very much enjoyed your post, however my understanding of quantum mechanics is minimal; and so I need to look further before I can truly understand. Is there any reading material you can refer one to regarding quantum mechanics and its relation with traditional metaphysics? I am primarily unfamiliar with quantum mechanics. I hope to receive a response.

\hfill

\texttt{Cologero on 2017-02-04 at 18:38 said:}

There is only one book that deals with quantum mechanics from the perspective of Tradition: \emph{The Quantum Enigma: Finding the Hidden Key}, by \textbf{Wolfgang Smith}.

\hfill

\texttt{Sigurd on 2017-02-04 at 19:27 said:}

Thanks for the response; and some further questions: is an understanding of Aristotelian metaphysics important to understanding Guenon's arguments in the Reign of Quantity? Would you consider Aristotle's metaphysics to be relevant today (including under a quantum framework)? How important is Thomism in this regard?

Thanks for the responses. This is also a good time to say I appreciate your work here. This site has been quite instrumental in the development of my understanding of Traditionalism, as well as placing them in their historical context vis-a-vis the letters of Evola and Guenon.

\hfill

\texttt{Cologero on 2017-02-04 at 19:45 said:}

Aristotle is important to understanding Guenon in the sense that it is more natural, as well as more accessible, to Westerners. For example, his doctrines of essence vs existence, the superiority of intellectual intuition, and the identification of knowledge with being, are part of Traditional metaphysics. Unfortunately, probably few professors in the West grasp the significance of those doctrines. Moreover, the Scholastics, like St Thomas Aquinas, developed those doctrines in a more Traditional direction. Look up Edward Feser for a basic introduction to Aristotle and Aquinas.

You can do quantum physics just fine with no knowledge of Aristotle, but that is not the point. Read the articles on Gornahoor about Wolfgang Smith for some ideas.

\hfill

\texttt{Sigurd on 2017-02-06 at 04:43 said:}

I see, thank you very much. One last distinction: would the divine creative act be likened to Aristotle's idea of `formal cause'? And would this phenomenon also be related to Evola's (and others) idea of involution? Apologies if the relevance of the questions are somewhat stretched, but I am trying to understand the essential creative act in a way that is congenial to my own readings. Again, thanks for the helpful responses.

\hfill

\texttt{Matt on 2017-02-07 at 19:34 said:}

Hi Sigurd,

I think the appropriate Aristotelian category to place the divine creative act in would be the efficient cause (also known as agent cause) rather than formal cause, which refers to the intrinsic nature (essential constitution) of a thing.

\hfill

\texttt{Sigurd on 2017-02-11 at 01:31 said:}

Great! Helpful response, Matt. Thank you for getting back to me. It is nice to find a website with such congenial spirits.

\end{sffamily}
\end{footnotesize}

